\honest{The Gift We Give To Tomorrow}{Eliezer Yudkowsky}{2008}

This post is a dialogue. ``I'' ask how natural selection, a cruel and mindless process, produced beings capable of love. ``You'' answer. I will call you the answerer. Selection is bloody and stupid: what counts is having more children than others, and a species can evolve to extinction if the winning genes play negative-sum games. \nb{The answerer knows the science and gets the last word, and the reader is addressed as the answerer. So the reader is on the winning side before any argument starts.}

My questions get harder. A mother's children share her genes. But mothers love adopted children, and love their children for themselves. We are adaptation-executers, not fitness-maximizers, says the answerer; through most of history no one knew genes existed. People befriend non-relatives; hunter-gatherers play iterated Prisoner's Dilemmas, and the most formidable person is often the one with most allies. Some die for their friends; we can tell true friends from fair-weather ones, and someone with many true friends is more formidable.

Then I ask about Gandhi, who turned the other cheek. We are political animals who argue about policy, says the answerer, and can argue ``What should be done?'' as a proposition. The answerer says Gandhi ``Believed certain complicated propositions about `What should be done?' and did them.'' I reply: ``That sounds like it could explain any possible human behavior.'' The answerer lists possible causes, a moral architecture able to argue abstract propositions, hardwired intuitions of fairness, duty, pain aversion and empathy, ``something like'' a preference for simple propositions, ``probably'' reused from our Occam prior, ``plus perhaps'' memetic selection, ending in ``You should not hurt people'' in full generality, and then says: ``Unless you think it was magic, it has to fit into the lawful causal development of the universe somehow.'' \nb{The questioner had not suggested magic.} I promise not to, and I do not ask about Gandhi again.

Then I marvel that hundreds of millions of years of evolution's death tournament produced mothers and fathers, steadfast friends and honourable enemies, true altruists, police officers and artists who sacrifice for their art. The answerer: if that surprises you, question your model, for ``Since the beginning, not one unusual thing has ever happened.'' A shadowy figure directing evolution would itself need to have evolved its love of love, and our loves bear the design signature of selection on the savanna.

Have I looked at this planet lately? Humans are not always nice. But we are far nicer than the process that let elephants starve when their teeth wore out and left dying gazelles unanaesthetized. A single twinge of empathy is more than evolution has. Beauty? Bach's \textsc{Little Fugue in G Minor} carries no tags of beauty in its sound waves. The only explicit measure of its beauty is in a human brain. The answerer says my question is circular. Evolution cannot quine its goal system. Making the first minds, its simple criterion shattered into a thousand values, including care for life and happiness, which evolution does not care about. I value beauty and altruism because I evolved to, so of course I find it marvelous that evolution produced them, and no coincidence needs explaining. Arguments for beauty and altruism ``from first principles'' are only different principles, which would not move a ghost of perfect emptiness either.

Supposing we tap into something beyond passes the recursive buck: why heed it more than our humanity, and how would following its orders change our responsibility? A standard of value outside human minds that happened to agree with ours would be ``far too much coincidence.'' \nb{Moral philosophers call this an evolutionary debunking argument. Sharon Street published a version in 2006, and philosophers have been replying to it since. Realists reply that tracking such a standard could itself have helped our ancestors survive, so the agreement need not be a coincidence. In the dialogue the argument wins by rhetorical question, and this reply is not considered.}

Next the answerer explains why flowers are beautiful. They ``evolved to attract bees---by imitating certain mating signs of bees.'' \nb{Flowers that attract insects by imitating their mating signals are known almost only among orchids; the known cases outside them are a South African daisy and an Asian iris. And many flowers are pollinated by other animals, or by wind.} Humans, the answerer goes on, like flowers because healthy flowers were a sign of fertile land. Before this story, the answerer asks whether I think ``science does not know the story.'' \nb{The part about humans is a speculative hypothesis, and the post gives no evidence for it.} Asked whether this explains the flower away, the answerer says: ``No, I explain it.'' When I call the creation of love ``a moral miracle,'' the answerer agrees, and says that if it still seems less wonderful to me, I ``have problems taking joy in the merely real.''

Love has to enter the universe somehow, as life did: trace your ancestry back far enough and you reach a replicator that arose by accident 3.85 billion years ago in some tidal pool. The answerer also states a law, supported only by that analogy: ``A complex pattern must be explained by a cause which is not already that complex pattern.'' \nb{Anyone who thinks love comes from a source that already loves is ruled out by the law, not by evidence.}

I end with a fairy tale. Our descendants will ask how they can love. Because we, who love, made them to. And how can we love? Because our parents, who loved, made us so. And where did it begin? They will be told that once ``there were intelligent beings who were not themselves intelligently designed,'' back ``when all of civilization was a single galaxy.'' \nb{The story, which begins with a ``perhaps,'' assumes a future of people designed by other people and a civilization spread beyond one galaxy. The post does not argue for either; they appear as background in a bedtime story.}
